% LaTeX Table Templates for SHAKTI-CHAIN Publication
% Include with: % LaTeX Table Templates for SHAKTI-CHAIN Publication
% Include with: % LaTeX Table Templates for SHAKTI-CHAIN Publication
% Include with: % LaTeX Table Templates for SHAKTI-CHAIN Publication
% Include with: \input{tables.tex}

% ============================================================
% HYPOTHESIS RESULTS TABLE
% ============================================================
% Usage: Use within document after defining data
%
% Example:
% \hypothesistable{Domain Name}{tab:domain1}{
%   H1.1 & Description & t-test & 2.45 & 0.012 & Supported \\
%   H1.2 & Description & t-test & 1.23 & 0.089 & Not Supp. \\
% }

\newcommand{\hypothesistable}[3]{
\begin{table}[htbp]
    \centering
    \caption{Hypothesis test results for #1}
    \label{#2}
    \small
    \begin{tabular}{@{}lp{2.5cm}lrrr@{}}
        \toprule
        \textbf{ID} & \textbf{Hypothesis} & \textbf{Test} & \textbf{Stat.} & \textbf{p-value} & \textbf{Decision} \\
        \midrule
        #3
        \bottomrule
    \end{tabular}
\end{table}
}

% ============================================================
% SUMMARY STATISTICS TABLE
% ============================================================
% Usage:
% \summarystats{Metric Name}{tab:summary}{
%   n & Mean & SD & Min & Q1 & Median & Q3 & Max \\
%   100 & 0.85 & 0.12 & 0.45 & 0.78 & 0.86 & 0.92 & 0.99 \\
% }

\newcommand{\summarystats}[3]{
\begin{table}[htbp]
    \centering
    \caption{Summary statistics for #1}
    \label{#2}
    \small
    \begin{tabular}{@{}rrrrrrrr@{}}
        \toprule
        \textbf{n} & \textbf{Mean} & \textbf{SD} & \textbf{Min} & \textbf{Q1} & \textbf{Median} & \textbf{Q3} & \textbf{Max} \\
        \midrule
        #3
        \bottomrule
    \end{tabular}
\end{table}
}

% ============================================================
% SYSTEM COMPARISON TABLE
% ============================================================
% Multi-system comparison with best values highlighted
%
% Usage:
% \comparisonTable{tab:comparison}{
%   Metric & SHAKTI & Fixed & Uniform & CDA \\
%   Efficiency & \textbf{0.94} & 0.78 & 0.82 & 0.89 \\
%   Fairness & \textbf{0.91} & 0.65 & 0.72 & 0.85 \\
% }

\newcommand{\comparisonTable}[2]{
\begin{table}[htbp]
    \centering
    \caption{System comparison across key metrics. Bold indicates best performance.}
    \label{#1}
    \begin{tabular}{@{}lcccc@{}}
        \toprule
        \textbf{Metric} & \textbf{SHAKTI} & \textbf{Fixed} & \textbf{Uniform} & \textbf{CDA} \\
        \midrule
        #2
        \bottomrule
    \end{tabular}
\end{table}
}

% ============================================================
% DOMAIN SUMMARY TABLE
% ============================================================
% Overview of results by domain

\newcommand{\domainsummary}[2]{
\begin{table}[htbp]
    \centering
    \caption{Summary of hypothesis test results by research domain}
    \label{#1}
    \begin{tabular}{@{}lcccr@{}}
        \toprule
        \textbf{Domain} & \textbf{Tested} & \textbf{Supported} & \textbf{Failed} & \textbf{Rate} \\
        \midrule
        #2
        \bottomrule
    \end{tabular}
\end{table}
}

% ============================================================
% EFFECT SIZE TABLE
% ============================================================
% Table for reporting effect sizes with confidence intervals

\newcommand{\effectsizetable}[2]{
\begin{table}[htbp]
    \centering
    \caption{Effect sizes with 95\% confidence intervals}
    \label{#1}
    \small
    \begin{tabular}{@{}llrrl@{}}
        \toprule
        \textbf{Hypothesis} & \textbf{Measure} & \textbf{Effect} & \textbf{95\% CI} & \textbf{Interpretation} \\
        \midrule
        #2
        \bottomrule
    \end{tabular}
\end{table}
}

% ============================================================
% POWER ANALYSIS TABLE
% ============================================================

\newcommand{\poweranalysis}[2]{
\begin{table}[htbp]
    \centering
    \caption{Statistical power analysis results}
    \label{#1}
    \small
    \begin{tabular}{@{}lrrrr@{}}
        \toprule
        \textbf{Test} & \textbf{n} & \textbf{Effect (d)} & \textbf{$\alpha$} & \textbf{Power} \\
        \midrule
        #2
        \bottomrule
    \end{tabular}
\end{table}
}

% ============================================================
% CRITICAL FAILURES TABLE
% ============================================================

\newcommand{\criticalfailures}[2]{
\begin{table}[htbp]
    \centering
    \caption{Critical hypothesis failures requiring attention}
    \label{#1}
    \small
    \begin{tabular}{@{}llp{4cm}l@{}}
        \toprule
        \textbf{ID} & \textbf{Domain} & \textbf{Issue} & \textbf{Priority} \\
        \midrule
        #2
        \bottomrule
    \end{tabular}
\end{table}
}

% ============================================================
% STYLING COMMANDS
% ============================================================

% Highlight significant p-values
\newcommand{\sigp}[1]{\textbf{#1}}

% Highlight supported hypotheses
\newcommand{\supported}{\textcolor{shaktigreen}{Supported}}
\newcommand{\notsupported}{\textcolor{shaktired}{Not Supp.}}

% Effect size interpretation
\newcommand{\smalleffect}{Small}
\newcommand{\mediumeffect}{\textbf{Medium}}
\newcommand{\largeeffect}{\textbf{\textcolor{shaktiblue}{Large}}}

% ============================================================
% LANDSCAPE TABLE ENVIRONMENT
% ============================================================
% For wide tables that need landscape orientation

\usepackage{pdflscape}

\newenvironment{landscapetable}{%
    \begin{landscape}
    \centering
}{%
    \end{landscape}
}

% ============================================================
% MULTI-PAGE TABLE SUPPORT
% ============================================================
\usepackage{longtable}

% Example usage:
% \begin{longtable}{@{}lp{3cm}lrrr@{}}
%     \caption{Full hypothesis results} \label{tab:full_results} \\
%     \toprule
%     ID & Hypothesis & Test & Stat. & p-value & Decision \\
%     \midrule
%     \endfirsthead
%     \multicolumn{6}{c}{\tablename\ \thetable{} -- continued} \\
%     \toprule
%     ID & Hypothesis & Test & Stat. & p-value & Decision \\
%     \midrule
%     \endhead
%     \midrule
%     \multicolumn{6}{r}{Continued on next page...} \\
%     \endfoot
%     \bottomrule
%     \endlastfoot
%     % Data rows here
% \end{longtable}


% ============================================================
% HYPOTHESIS RESULTS TABLE
% ============================================================
% Usage: Use within document after defining data
%
% Example:
% \hypothesistable{Domain Name}{tab:domain1}{
%   H1.1 & Description & t-test & 2.45 & 0.012 & Supported \\
%   H1.2 & Description & t-test & 1.23 & 0.089 & Not Supp. \\
% }

\newcommand{\hypothesistable}[3]{
\begin{table}[htbp]
    \centering
    \caption{Hypothesis test results for #1}
    \label{#2}
    \small
    \begin{tabular}{@{}lp{2.5cm}lrrr@{}}
        \toprule
        \textbf{ID} & \textbf{Hypothesis} & \textbf{Test} & \textbf{Stat.} & \textbf{p-value} & \textbf{Decision} \\
        \midrule
        #3
        \bottomrule
    \end{tabular}
\end{table}
}

% ============================================================
% SUMMARY STATISTICS TABLE
% ============================================================
% Usage:
% \summarystats{Metric Name}{tab:summary}{
%   n & Mean & SD & Min & Q1 & Median & Q3 & Max \\
%   100 & 0.85 & 0.12 & 0.45 & 0.78 & 0.86 & 0.92 & 0.99 \\
% }

\newcommand{\summarystats}[3]{
\begin{table}[htbp]
    \centering
    \caption{Summary statistics for #1}
    \label{#2}
    \small
    \begin{tabular}{@{}rrrrrrrr@{}}
        \toprule
        \textbf{n} & \textbf{Mean} & \textbf{SD} & \textbf{Min} & \textbf{Q1} & \textbf{Median} & \textbf{Q3} & \textbf{Max} \\
        \midrule
        #3
        \bottomrule
    \end{tabular}
\end{table}
}

% ============================================================
% SYSTEM COMPARISON TABLE
% ============================================================
% Multi-system comparison with best values highlighted
%
% Usage:
% \comparisonTable{tab:comparison}{
%   Metric & SHAKTI & Fixed & Uniform & CDA \\
%   Efficiency & \textbf{0.94} & 0.78 & 0.82 & 0.89 \\
%   Fairness & \textbf{0.91} & 0.65 & 0.72 & 0.85 \\
% }

\newcommand{\comparisonTable}[2]{
\begin{table}[htbp]
    \centering
    \caption{System comparison across key metrics. Bold indicates best performance.}
    \label{#1}
    \begin{tabular}{@{}lcccc@{}}
        \toprule
        \textbf{Metric} & \textbf{SHAKTI} & \textbf{Fixed} & \textbf{Uniform} & \textbf{CDA} \\
        \midrule
        #2
        \bottomrule
    \end{tabular}
\end{table}
}

% ============================================================
% DOMAIN SUMMARY TABLE
% ============================================================
% Overview of results by domain

\newcommand{\domainsummary}[2]{
\begin{table}[htbp]
    \centering
    \caption{Summary of hypothesis test results by research domain}
    \label{#1}
    \begin{tabular}{@{}lcccr@{}}
        \toprule
        \textbf{Domain} & \textbf{Tested} & \textbf{Supported} & \textbf{Failed} & \textbf{Rate} \\
        \midrule
        #2
        \bottomrule
    \end{tabular}
\end{table}
}

% ============================================================
% EFFECT SIZE TABLE
% ============================================================
% Table for reporting effect sizes with confidence intervals

\newcommand{\effectsizetable}[2]{
\begin{table}[htbp]
    \centering
    \caption{Effect sizes with 95\% confidence intervals}
    \label{#1}
    \small
    \begin{tabular}{@{}llrrl@{}}
        \toprule
        \textbf{Hypothesis} & \textbf{Measure} & \textbf{Effect} & \textbf{95\% CI} & \textbf{Interpretation} \\
        \midrule
        #2
        \bottomrule
    \end{tabular}
\end{table}
}

% ============================================================
% POWER ANALYSIS TABLE
% ============================================================

\newcommand{\poweranalysis}[2]{
\begin{table}[htbp]
    \centering
    \caption{Statistical power analysis results}
    \label{#1}
    \small
    \begin{tabular}{@{}lrrrr@{}}
        \toprule
        \textbf{Test} & \textbf{n} & \textbf{Effect (d)} & \textbf{$\alpha$} & \textbf{Power} \\
        \midrule
        #2
        \bottomrule
    \end{tabular}
\end{table}
}

% ============================================================
% CRITICAL FAILURES TABLE
% ============================================================

\newcommand{\criticalfailures}[2]{
\begin{table}[htbp]
    \centering
    \caption{Critical hypothesis failures requiring attention}
    \label{#1}
    \small
    \begin{tabular}{@{}llp{4cm}l@{}}
        \toprule
        \textbf{ID} & \textbf{Domain} & \textbf{Issue} & \textbf{Priority} \\
        \midrule
        #2
        \bottomrule
    \end{tabular}
\end{table}
}

% ============================================================
% STYLING COMMANDS
% ============================================================

% Highlight significant p-values
\newcommand{\sigp}[1]{\textbf{#1}}

% Highlight supported hypotheses
\newcommand{\supported}{\textcolor{shaktigreen}{Supported}}
\newcommand{\notsupported}{\textcolor{shaktired}{Not Supp.}}

% Effect size interpretation
\newcommand{\smalleffect}{Small}
\newcommand{\mediumeffect}{\textbf{Medium}}
\newcommand{\largeeffect}{\textbf{\textcolor{shaktiblue}{Large}}}

% ============================================================
% LANDSCAPE TABLE ENVIRONMENT
% ============================================================
% For wide tables that need landscape orientation

\usepackage{pdflscape}

\newenvironment{landscapetable}{%
    \begin{landscape}
    \centering
}{%
    \end{landscape}
}

% ============================================================
% MULTI-PAGE TABLE SUPPORT
% ============================================================
\usepackage{longtable}

% Example usage:
% \begin{longtable}{@{}lp{3cm}lrrr@{}}
%     \caption{Full hypothesis results} \label{tab:full_results} \\
%     \toprule
%     ID & Hypothesis & Test & Stat. & p-value & Decision \\
%     \midrule
%     \endfirsthead
%     \multicolumn{6}{c}{\tablename\ \thetable{} -- continued} \\
%     \toprule
%     ID & Hypothesis & Test & Stat. & p-value & Decision \\
%     \midrule
%     \endhead
%     \midrule
%     \multicolumn{6}{r}{Continued on next page...} \\
%     \endfoot
%     \bottomrule
%     \endlastfoot
%     % Data rows here
% \end{longtable}


% ============================================================
% HYPOTHESIS RESULTS TABLE
% ============================================================
% Usage: Use within document after defining data
%
% Example:
% \hypothesistable{Domain Name}{tab:domain1}{
%   H1.1 & Description & t-test & 2.45 & 0.012 & Supported \\
%   H1.2 & Description & t-test & 1.23 & 0.089 & Not Supp. \\
% }

\newcommand{\hypothesistable}[3]{
\begin{table}[htbp]
    \centering
    \caption{Hypothesis test results for #1}
    \label{#2}
    \small
    \begin{tabular}{@{}lp{2.5cm}lrrr@{}}
        \toprule
        \textbf{ID} & \textbf{Hypothesis} & \textbf{Test} & \textbf{Stat.} & \textbf{p-value} & \textbf{Decision} \\
        \midrule
        #3
        \bottomrule
    \end{tabular}
\end{table}
}

% ============================================================
% SUMMARY STATISTICS TABLE
% ============================================================
% Usage:
% \summarystats{Metric Name}{tab:summary}{
%   n & Mean & SD & Min & Q1 & Median & Q3 & Max \\
%   100 & 0.85 & 0.12 & 0.45 & 0.78 & 0.86 & 0.92 & 0.99 \\
% }

\newcommand{\summarystats}[3]{
\begin{table}[htbp]
    \centering
    \caption{Summary statistics for #1}
    \label{#2}
    \small
    \begin{tabular}{@{}rrrrrrrr@{}}
        \toprule
        \textbf{n} & \textbf{Mean} & \textbf{SD} & \textbf{Min} & \textbf{Q1} & \textbf{Median} & \textbf{Q3} & \textbf{Max} \\
        \midrule
        #3
        \bottomrule
    \end{tabular}
\end{table}
}

% ============================================================
% SYSTEM COMPARISON TABLE
% ============================================================
% Multi-system comparison with best values highlighted
%
% Usage:
% \comparisonTable{tab:comparison}{
%   Metric & SHAKTI & Fixed & Uniform & CDA \\
%   Efficiency & \textbf{0.94} & 0.78 & 0.82 & 0.89 \\
%   Fairness & \textbf{0.91} & 0.65 & 0.72 & 0.85 \\
% }

\newcommand{\comparisonTable}[2]{
\begin{table}[htbp]
    \centering
    \caption{System comparison across key metrics. Bold indicates best performance.}
    \label{#1}
    \begin{tabular}{@{}lcccc@{}}
        \toprule
        \textbf{Metric} & \textbf{SHAKTI} & \textbf{Fixed} & \textbf{Uniform} & \textbf{CDA} \\
        \midrule
        #2
        \bottomrule
    \end{tabular}
\end{table}
}

% ============================================================
% DOMAIN SUMMARY TABLE
% ============================================================
% Overview of results by domain

\newcommand{\domainsummary}[2]{
\begin{table}[htbp]
    \centering
    \caption{Summary of hypothesis test results by research domain}
    \label{#1}
    \begin{tabular}{@{}lcccr@{}}
        \toprule
        \textbf{Domain} & \textbf{Tested} & \textbf{Supported} & \textbf{Failed} & \textbf{Rate} \\
        \midrule
        #2
        \bottomrule
    \end{tabular}
\end{table}
}

% ============================================================
% EFFECT SIZE TABLE
% ============================================================
% Table for reporting effect sizes with confidence intervals

\newcommand{\effectsizetable}[2]{
\begin{table}[htbp]
    \centering
    \caption{Effect sizes with 95\% confidence intervals}
    \label{#1}
    \small
    \begin{tabular}{@{}llrrl@{}}
        \toprule
        \textbf{Hypothesis} & \textbf{Measure} & \textbf{Effect} & \textbf{95\% CI} & \textbf{Interpretation} \\
        \midrule
        #2
        \bottomrule
    \end{tabular}
\end{table}
}

% ============================================================
% POWER ANALYSIS TABLE
% ============================================================

\newcommand{\poweranalysis}[2]{
\begin{table}[htbp]
    \centering
    \caption{Statistical power analysis results}
    \label{#1}
    \small
    \begin{tabular}{@{}lrrrr@{}}
        \toprule
        \textbf{Test} & \textbf{n} & \textbf{Effect (d)} & \textbf{$\alpha$} & \textbf{Power} \\
        \midrule
        #2
        \bottomrule
    \end{tabular}
\end{table}
}

% ============================================================
% CRITICAL FAILURES TABLE
% ============================================================

\newcommand{\criticalfailures}[2]{
\begin{table}[htbp]
    \centering
    \caption{Critical hypothesis failures requiring attention}
    \label{#1}
    \small
    \begin{tabular}{@{}llp{4cm}l@{}}
        \toprule
        \textbf{ID} & \textbf{Domain} & \textbf{Issue} & \textbf{Priority} \\
        \midrule
        #2
        \bottomrule
    \end{tabular}
\end{table}
}

% ============================================================
% STYLING COMMANDS
% ============================================================

% Highlight significant p-values
\newcommand{\sigp}[1]{\textbf{#1}}

% Highlight supported hypotheses
\newcommand{\supported}{\textcolor{shaktigreen}{Supported}}
\newcommand{\notsupported}{\textcolor{shaktired}{Not Supp.}}

% Effect size interpretation
\newcommand{\smalleffect}{Small}
\newcommand{\mediumeffect}{\textbf{Medium}}
\newcommand{\largeeffect}{\textbf{\textcolor{shaktiblue}{Large}}}

% ============================================================
% LANDSCAPE TABLE ENVIRONMENT
% ============================================================
% For wide tables that need landscape orientation

\usepackage{pdflscape}

\newenvironment{landscapetable}{%
    \begin{landscape}
    \centering
}{%
    \end{landscape}
}

% ============================================================
% MULTI-PAGE TABLE SUPPORT
% ============================================================
\usepackage{longtable}

% Example usage:
% \begin{longtable}{@{}lp{3cm}lrrr@{}}
%     \caption{Full hypothesis results} \label{tab:full_results} \\
%     \toprule
%     ID & Hypothesis & Test & Stat. & p-value & Decision \\
%     \midrule
%     \endfirsthead
%     \multicolumn{6}{c}{\tablename\ \thetable{} -- continued} \\
%     \toprule
%     ID & Hypothesis & Test & Stat. & p-value & Decision \\
%     \midrule
%     \endhead
%     \midrule
%     \multicolumn{6}{r}{Continued on next page...} \\
%     \endfoot
%     \bottomrule
%     \endlastfoot
%     % Data rows here
% \end{longtable}


% ============================================================
% HYPOTHESIS RESULTS TABLE
% ============================================================
% Usage: Use within document after defining data
%
% Example:
% \hypothesistable{Domain Name}{tab:domain1}{
%   H1.1 & Description & t-test & 2.45 & 0.012 & Supported \\
%   H1.2 & Description & t-test & 1.23 & 0.089 & Not Supp. \\
% }

\newcommand{\hypothesistable}[3]{
\begin{table}[htbp]
    \centering
    \caption{Hypothesis test results for #1}
    \label{#2}
    \small
    \begin{tabular}{@{}lp{2.5cm}lrrr@{}}
        \toprule
        \textbf{ID} & \textbf{Hypothesis} & \textbf{Test} & \textbf{Stat.} & \textbf{p-value} & \textbf{Decision} \\
        \midrule
        #3
        \bottomrule
    \end{tabular}
\end{table}
}

% ============================================================
% SUMMARY STATISTICS TABLE
% ============================================================
% Usage:
% \summarystats{Metric Name}{tab:summary}{
%   n & Mean & SD & Min & Q1 & Median & Q3 & Max \\
%   100 & 0.85 & 0.12 & 0.45 & 0.78 & 0.86 & 0.92 & 0.99 \\
% }

\newcommand{\summarystats}[3]{
\begin{table}[htbp]
    \centering
    \caption{Summary statistics for #1}
    \label{#2}
    \small
    \begin{tabular}{@{}rrrrrrrr@{}}
        \toprule
        \textbf{n} & \textbf{Mean} & \textbf{SD} & \textbf{Min} & \textbf{Q1} & \textbf{Median} & \textbf{Q3} & \textbf{Max} \\
        \midrule
        #3
        \bottomrule
    \end{tabular}
\end{table}
}

% ============================================================
% SYSTEM COMPARISON TABLE
% ============================================================
% Multi-system comparison with best values highlighted
%
% Usage:
% \comparisonTable{tab:comparison}{
%   Metric & SHAKTI & Fixed & Uniform & CDA \\
%   Efficiency & \textbf{0.94} & 0.78 & 0.82 & 0.89 \\
%   Fairness & \textbf{0.91} & 0.65 & 0.72 & 0.85 \\
% }

\newcommand{\comparisonTable}[2]{
\begin{table}[htbp]
    \centering
    \caption{System comparison across key metrics. Bold indicates best performance.}
    \label{#1}
    \begin{tabular}{@{}lcccc@{}}
        \toprule
        \textbf{Metric} & \textbf{SHAKTI} & \textbf{Fixed} & \textbf{Uniform} & \textbf{CDA} \\
        \midrule
        #2
        \bottomrule
    \end{tabular}
\end{table}
}

% ============================================================
% DOMAIN SUMMARY TABLE
% ============================================================
% Overview of results by domain

\newcommand{\domainsummary}[2]{
\begin{table}[htbp]
    \centering
    \caption{Summary of hypothesis test results by research domain}
    \label{#1}
    \begin{tabular}{@{}lcccr@{}}
        \toprule
        \textbf{Domain} & \textbf{Tested} & \textbf{Supported} & \textbf{Failed} & \textbf{Rate} \\
        \midrule
        #2
        \bottomrule
    \end{tabular}
\end{table}
}

% ============================================================
% EFFECT SIZE TABLE
% ============================================================
% Table for reporting effect sizes with confidence intervals

\newcommand{\effectsizetable}[2]{
\begin{table}[htbp]
    \centering
    \caption{Effect sizes with 95\% confidence intervals}
    \label{#1}
    \small
    \begin{tabular}{@{}llrrl@{}}
        \toprule
        \textbf{Hypothesis} & \textbf{Measure} & \textbf{Effect} & \textbf{95\% CI} & \textbf{Interpretation} \\
        \midrule
        #2
        \bottomrule
    \end{tabular}
\end{table}
}

% ============================================================
% POWER ANALYSIS TABLE
% ============================================================

\newcommand{\poweranalysis}[2]{
\begin{table}[htbp]
    \centering
    \caption{Statistical power analysis results}
    \label{#1}
    \small
    \begin{tabular}{@{}lrrrr@{}}
        \toprule
        \textbf{Test} & \textbf{n} & \textbf{Effect (d)} & \textbf{$\alpha$} & \textbf{Power} \\
        \midrule
        #2
        \bottomrule
    \end{tabular}
\end{table}
}

% ============================================================
% CRITICAL FAILURES TABLE
% ============================================================

\newcommand{\criticalfailures}[2]{
\begin{table}[htbp]
    \centering
    \caption{Critical hypothesis failures requiring attention}
    \label{#1}
    \small
    \begin{tabular}{@{}llp{4cm}l@{}}
        \toprule
        \textbf{ID} & \textbf{Domain} & \textbf{Issue} & \textbf{Priority} \\
        \midrule
        #2
        \bottomrule
    \end{tabular}
\end{table}
}

% ============================================================
% STYLING COMMANDS
% ============================================================

% Highlight significant p-values
\newcommand{\sigp}[1]{\textbf{#1}}

% Highlight supported hypotheses
\newcommand{\supported}{\textcolor{shaktigreen}{Supported}}
\newcommand{\notsupported}{\textcolor{shaktired}{Not Supp.}}

% Effect size interpretation
\newcommand{\smalleffect}{Small}
\newcommand{\mediumeffect}{\textbf{Medium}}
\newcommand{\largeeffect}{\textbf{\textcolor{shaktiblue}{Large}}}

% ============================================================
% LANDSCAPE TABLE ENVIRONMENT
% ============================================================
% For wide tables that need landscape orientation

\usepackage{pdflscape}

\newenvironment{landscapetable}{%
    \begin{landscape}
    \centering
}{%
    \end{landscape}
}

% ============================================================
% MULTI-PAGE TABLE SUPPORT
% ============================================================
\usepackage{longtable}

% Example usage:
% \begin{longtable}{@{}lp{3cm}lrrr@{}}
%     \caption{Full hypothesis results} \label{tab:full_results} \\
%     \toprule
%     ID & Hypothesis & Test & Stat. & p-value & Decision \\
%     \midrule
%     \endfirsthead
%     \multicolumn{6}{c}{\tablename\ \thetable{} -- continued} \\
%     \toprule
%     ID & Hypothesis & Test & Stat. & p-value & Decision \\
%     \midrule
%     \endhead
%     \midrule
%     \multicolumn{6}{r}{Continued on next page...} \\
%     \endfoot
%     \bottomrule
%     \endlastfoot
%     % Data rows here
% \end{longtable}
